\section{Conclusion and Discussion}
\label{sec:conclusion}

In this paper, we introduced \textbf{FlatQ-Merge}, an exhaustive empirical study and framework investigating the role of expert loss landscape geometry in low-precision model merging. Through systematic experiments across 3 independent seeds on a Vision Transformer backbone spanning MNIST, FashionMNIST, CIFAR-10, and SVHN, we demonstrated that promoting wider and flatter local minima during expert pre-training directly enhances resilience to quantization noise. Crucially, this synergy is highly precision-dependent: under standard 8-bit quantization, SGD-trained experts are highly robust and do not benefit significantly from flatness, whereas under aggressive 4-bit quantization, our optimal configuration ($\rho=0.05$) achieved a substantial \textbf{+7.44\%} absolute accuracy improvement over standard SGD-trained experts. Additionally, we identified a critical non-linear \emph{Over-Perturbation Threshold} at $\rho \ge 0.1$, where excessive perturbation destabilizes training and degrades task-vector quality.

\subsection{Limitations and Scope}
We explicitly outline the limitations and boundaries of our evaluation:
\begin{enumerate}
  \item \textbf{Data and Scale Constraints and Absolute Accuracy Gap:} To make our multi-axial grid sweeps computationally tractable, we restricted pre-training to 512 images per task on a compact \texttt{vit\_tiny\_patch16\_224} backbone. While this simulates extreme low-data edge deployment, absolute test accuracies are degraded compared to full-scale, converged models (e.g., individual unquantized experts achieve ~64.28% on this budget, while fully converged individual models on the entire dataset achieve ~95%). To bridge this absolute performance gap, future practitioners can scale the pre-training dataset size (e.g., using 5,000 or more pre-training samples) or utilize slightly larger backbones (such as ViT-Small or ViT-Base). Because the relative accuracy gains of +7.44% under 4-bit quantization are statistically significant and driven by the physical properties of loss landscape geometry, we anticipate that these gains will be fully preserved or even amplified at larger scales. Furthermore, because modern Large Language Models (LLMs) are notoriously sensitive to quantization outliers, we anticipate that our flatness-robustness synergy will generalize and potentially strengthen when scaling to billion-parameter autoregressive models. Specifically, modern LLM post-training quantization (PTQ) techniques (such as AWQ \cite{Lin2023AWQ} or GPTQ \cite{Frantar2022GPTQ}) face severe bottlenecks from activation outliers and weight-rounding errors. By incorporating SAM pre-training or landscape-flattening objectives into the instruction-tuning phase of these models, we can structurally suppress outlier channels and smooth the parameter manifold, allowing downstream PTQ to achieve nearly lossless low-bit compression without requiring complex test-time search.
  \item \textbf{Low-Latency Optimization Window:} Test-time adaptation was restricted to 40 steps to maintain negligible deployment latency. Consequently, the coefficients reside in a highly complex local saddle-point or flat plateau region rather than an asymptotic, perfectly stationary point. In a piecewise-constant quantized loss landscape, this local geometry combined with discretization noise leads to small, oscillating expected entropy changes under random perturbations.
  \item \textbf{Task Incongruence, Parameter Interference, and Domain Alignment:} Merging highly divergent datasets (e.g., MNIST and CIFAR-10) onto a compact backbone introduces extreme task incongruence, explaining the low absolute multi-task accuracies in our multi-task sandbox. However, real-world model merging typically occurs on highly related domains or standard domain-merging benchmarks (such as DomainNet or Office-Home, where experts are fine-tuned on different sub-domains of the same core category), where parameter interference is naturally and significantly reduced and absolute performance is preserved. In highly aligned domains, the task vectors are largely correlated, sharing similar sub-dimensional manifolds of the weight space. Under such conditions, we hypothesize that the flatness-robustness synergy will become even more pronounced. Because flat pre-training aligns the weights of different experts near a common wide valley, it minimizes the destructive interference of highly correlated task vectors when they are fused and quantized, preserving high absolute multi-domain performance.
  \item \textbf{Weight-Only Quantization:} Our framework focuses on weight-only post-training quantization (W8A32 and W4A32). While weight compression represents the primary memory saving on edge devices, edge deployment often requires joint weight-activation quantization (e.g., W8A8 or W4A4) for integer-only arithmetic. Future work will extend FlatQ-Merge to joint quantization by incorporating activation scaling factors or estimators like SmoothQuant. Importantly, we expect SAM-induced expert flatness to actively suppress activation outliers. By penalizing large weight gradients and smoothing local parameters, flat experts exhibit smaller weight magnitudes and spectral norms. Mathematically, bounding the spectral norm $\|W\|_2$ or the Frobenius norm $\|W\|_F$ of a layer's weight matrix directly restricts its Lipschitz constant, since the Lipschitz constant of a linear or convolutional layer is given by its spectral norm: $\text{Lip}(f) = \|W\|_2$. Because the maximum activation magnitude of a layer $y = f(x)$ is bounded by the layer's Lipschitz constant times the input range ($\|f(x)\| \le \text{Lip}(f) \|x\|$), restricting weight norms directly bounds the propagation of extreme activation spikes and prevents the emergence of massive activation outliers across successive transformer layers. This smoother activation distribution would inherently mitigate the representation clipping and rounding noise that typically plagues low-bit activation PTQ, representing a high-potential synergy.
  \item \textbf{Generalization Across Architectures (ViTs vs. CNNs):} Our empirical evaluation is conducted on a Vision Transformer (\texttt{vit\_tiny\_patch16\_224}) backbone. Vision Transformers are known to lack the strong translation-invariance and inductive biases of Convolutional Neural Networks (CNNs), resulting in a highly complex and non-convex loss landscape that is exceptionally sensitive to quantization rounding noise and training trajectories. While ViTs serve as an ideal, challenging empirical "sandbox" to stress-test our framework, the flatness-robustness synergy is fundamentally architecture-independent and rooted in second-order optimization theory (as derived in Section~\ref{sec:method}). The projection relation $H_{\Lambda} = T^T H_{\theta} T$ holds for any parameterized layer group, whether transformer attention projection or convolutional filters. We discuss the feasibility of scaling this framework to a ResNet-18 or other standard CNN backbones, where we anticipate that the strong inductive biases of CNNs will result in even higher absolute accuracies under both 8-bit and 4-bit precision, while fully preserving the +7.44\% absolute improvement of flat experts over standard SGD experts.
\end{enumerate}

\subsection{Future Directions}
Based on our empirical observations, we identify three high-priority future directions:
\begin{enumerate}
  \item \textbf{Layer-wise and Block-specific Flatness Control:} Currently, a uniform SAM radius $\rho$ is applied across all layers. However, model merging sensitivity is highly non-uniform across the network, with attention projection layers and deeper MLP blocks exhibiting the highest parameter variance and vulnerability to interference. Future work will explore dynamic layer-wise flatness schedules:
  \begin{equation}
    \rho_l = \rho_{\text{base}} \cdot \gamma(l)
  \end{equation}
  where $\gamma(l)$ scales the perturbation radius based on layer depth $l$ or layer-wise Fisher information. By applying a larger radius $\rho_l$ to highly sensitive multi-head attention blocks and a smaller radius to robust early projection layers, we can optimize the flatness-robustness trade-off and potentially bypass the Over-Perturbation Threshold ($\rho \ge 0.1$) on a global level.
  \item \textbf{Joint Weight-Activation Flatness:} Extending our framework to optimize flatness directly with respect to both weights and activations during pre-training. By incorporating activation-smoothing regularization (e.g., minimizing activation variance or bounding local Lipschitz constants), we can ensure that both weights and activation states are highly robust under joint low-bit (e.g., W4A4 or W8A8) quantization, facilitating deployment on integer-only hardware.
  \item \textbf{Scaling to LLMs and Autoregressive Models:} Applying flatness-inducing pre-training or instruction-tuning (e.g., via SAM or adversarial training) to large generative models (such as LLaMA or Mistral). Suppressing outlier channels at the source through flatness control holds immense potential to unlock painless, lossless 4-bit model merging and low-bit weight quantization in full-scale LLM deployments.
\end{enumerate}

By bridging the gap between loss landscape geometry and low-precision model merging, FlatQ-Merge provides a solid foundation and highly reproducible blueprint for deploying robust, parameter-fused, and highly compressed models in real-world scenarios.
